% CPSC 462 Computer Networks — course-specific macros ONLY.
%
% Inputted AFTER docs/latex/coursework_preamble.tex, which carries everything
% shared (answerbox, trap, problemheader, \srcref, \extref, page style).
% Put nothing general here: if another course would want it, it belongs in the
% shared preamble.
%
% ONE file for the whole course, at course level rather than under hw/. The
% syllabus grades homework, quizzes AND two exams, so hw/, qz/ and exam/ are
% siblings; a macros file inside any one of them would either be unreachable
% from the other two or get copied and drift.
%
% Doctrine: docs/directives/coursework-solutions.md

\renewcommand{\coursename}{CPSC 462}

% --- packet delay -----------------------------------------------------------
% The course's own notation, copied off the slide rather than invented:
%   dnodal = dproc + dqueue + dtrans + dprop     Introduction, PDF p50, slide 1-50
%   dtrans = L/R,  dprop = d/s                   Introduction, PDF p51, slide 1-51
% The slide sets an italic d with upright subscripts, which is what these give.
\newcommand{\dnodal}{d_{\mathrm{nodal}}}
\newcommand{\dproc}{d_{\mathrm{proc}}}
\newcommand{\dqueue}{d_{\mathrm{queue}}}
\newcommand{\dtrans}{d_{\mathrm{trans}}}
\newcommand{\dprop}{d_{\mathrm{prop}}}

% Round-trip time. Upright, because the slides treat it as a name and use it as
% a coefficient -- "non-persistent HTTP response time = 2RTT + file
% transmission time".                       Application Layer, PDF p29, slide 2-29
\newcommand{\RTT}{\mathrm{RTT}}

% --- link rate --------------------------------------------------------------
% The decks write "R = 100 Mb/s" (Introduction, PDF p33, slide 1-33). siunitx is
% loaded by the shared preamble, but \unit{\mega\bit\per\second} renders as
% "Mbit s^{-1}", which is not what this course writes. Hence one shorthand.
\newcommand{\Mbps}{\,\mathrm{Mb/s}}

% --- protocol / packet header field -----------------------------------------
% \field{Dest Port}, \field{TTL}, \field{http.request.method}. A packet-analysis
% course names fields and display filters constantly, and they must not be
% mistaken for prose or for a variable.
\newcommand{\field}[1]{\texttt{#1}}

% Kept deliberately small. Add a macro when a SECOND problem needs the same
% notation -- not in anticipation.
%
% Removed 2026-09-05 for want of a user (see reference_docs/findings.md):
%   \kbps  "kb/s" and "kbps" appear nowhere in the course notes. The nearest
%          hit is "L = 10 Kbits", which is a quantity of bits, not a rate.
%   \ms    a generic unit shorthand, not networks notation, and siunitx already
%          renders \unit{\milli\second} as "ms" correctly.
